\chapter{Introdução}


Historicamente, o conceito de Habilidade Social (HS) passou por diversas transformações e recebeu diferentes denominações, como ``comportamento assertivo'', ``competência interpessoal'' e ``efetividade pessoal'' \cite{Wolpe1968} \cite{Lazarus1971}. Os primeiros estudos empíricos datam do início do século XX, com observações sobre o comportamento social de crianças, mas foi a partir da década de 1970 que o treinamento em habilidades sociais (THS) ganhou destaque como técnica terapêutica, sendo incorporado em intervenções clínicas, escolares e organizacionais \cite{caballo2003manual}. Desde então, o THS tem sido utilizado como estratégia eficaz para reduzir \textit{déficits} de interação social e promover o bem-estar psicológico.

Sendo assim, de acordo com \citeautor{caballo2003manual}, as HS referem-se a um conjunto de comportamentos aprendidos que envolvem a expressão adequada de sentimentos, desejos, opiniões e direitos em contextos interpessoais, de modo que respeitem os outros e maximizem os resultados positivos da interação, e segundo o próprio autor ``o comportamento socialmente hábil é esse conjunto de comportamentos ... desse indivíduo de modo adequado à situação, respeitando esses comportamentos nos demais, e que geralmente resolve os problemas imediatos da situação enquanto minimiza a probabilidade de futuros problemas'' \citeautor{Caballo1996}.

A história dos videogames remonta à década de 1950, com os primeiros experimentos em simulações interativas realizadas em laboratórios de pesquisa. Um marco inicial foi o jogo ``Tennis for Two'' desenvolvido por William Higinbotham em 1958, seguido pelo icônico ``Pong'' da Atari na década de 1970, que ajudou a consolidar o setor de entretenimento eletrônico doméstico \cite{kent2001}. Esses jogos rudimentares pavimentaram o caminho para a popularização de \textit{consoles} como o Atari 2600, que transformou os videogames em um fenômeno cultural global.

Ao longo das décadas, os videogames evoluíram não apenas tecnicamente, mas também como forma de expressão cultural e artística. Autores como \citeautor{donovan2010} destacam que os jogos eletrônicos passaram a refletir narrativas complexas e a integrar elementos socioculturais diversos. A preservação e análise dessa história tornaram-se temas relevantes na academia, com iniciativas voltadas à memória e à documentação do desenvolvimento da indústria, como discutido por \citeautor{wolf2012} em sua coletânea sobre os primeiros anos dos jogos digitais.

Diversos estudos sugerem que jogar videogames pode trazer benefícios cognitivos, emocionais e sociais significativos. O artigo de  \citeautor{granic2014} destaca que esses jogos têm impacto positivo em quatro domínios principais: funções cognitivas (como atenção seletiva), motivação, regulação emocional e habilidades sociais. Além disso, pesquisas mostram que jogos casuais melhoram a memória e a agilidade mental entre adultos, com efeitos diferenciados por faixa etária — jovens percebem ganhos em memória e raciocínio, enquanto adultos mais velhos relatam melhoras em habilidades visuoespaciais e tempo de reação \cite{whitbourne2013}.

E entre os estudos sobre jogos, destacam-se os jogos \textit{multiplayer}. Esses que são especialmente online, e são ambientes interativos que envolvem múltiplos jogadores em tempo real, em modos cooperativos ou competitivos. São reconhecidos por promover um forte senso de comunidade e apoio social, inclusive em contextos como a pandemia de COVID-19, onde mantiveram conexões com amigos mesmo durante o isolamento \cite{cmentowski2020}. Estudos com jogos famosos como \textbf{Halo: Reach}, analisaram como que amizades dentro do jogo não apenas melhoram o desempenho individual e coletivo, mas também reforçam comportamentos pró-sociais e o engajamento contínuo dos jogadores \cite{mason2012}.

Desta forma, o presente trabalho tem como objetivo geral \textbf{investigar se jogos multiplayer podem ser utilizados como ferramentas de treinamento de habilidades sociais}.

\section{Estrutura do Trabalho}
\todo[inline]{colocar o que cada capítulo apresenta, resumidamente}